\setcounter{definition}{0}
\setcounter{theorem}{0}
\subsection{Bevezetés}
\subsubsection{Maradékos osztás}
\begin{theorem}
    Tetszőleges $a, b \neq 0$ egész számokhoz
    egyértelműen léteznek $q, r$ egészek, hogy
    \begin{equation*}
        a = b \cdot q + r \quad 0 \leq r < |b|
    \end{equation*}
\end{theorem}
\begin{remark*}
    Jelölése: $r = a \bmod b$
\end{remark*}
\begin{proof}
    \begin{equation*}
        a < b \implies a = b \cdot 0 + a \implies q = 0 \land r = a
    \end{equation*}
    \begin{multline*}
        a \geq b \implies \exists q', r': a - b = b \cdot q' + r' \implies \\
        a = b \cdot (q' + 1) + r' \implies q = q' + 1 \land r = r'
    \end{multline*}
    \begin{gather*}
        a = b \cdot q + r = b \cdot q' + r' \quad (r' \geq r) \implies
        b \cdot (q - q') = r' - r \\
        q \neq q' \implies |q - q'| > 1 \implies |b \cdot (q - q')| > b \implies
        |r' - r| = r' - r > b
        \tag{Ellentmondás, mert $r, r' < |b|$}
    \end{gather*}
\end{proof}
\subsubsection{Legnagyobb közös osztó}
\begin{definition}
    Legyenek $a, b \in \mathbb{Z}$.
    A $d$ nemnegatív egész az $a$ és $b$ legnagyobb közös osztója, ha
    \begin{itemize}
        \item $d \mid a$ és $d \mid b$
        \item minden $k \in \mathbb{Z}$ esetén, ha $k \mid a$ és $k \mid b$,
        akkor $k \mid d$
    \end{itemize}
    $a = b = 0$ esetén $d = 0$.
\end{definition}
\begin{remark*}
    Jelölése: $d = (a, b) = \text{lnko}(a, b) = \gcd(a, b)$
\end{remark*}
\begin{theorem}
    Minden $a, b \in \mathbb{Z}$ esetén
    létezik az $(a, b)$ legnagyobb közös osztó.
\end{theorem}
\begin{proof}
    \begin{equation*}
        (a, b) = (0, 0) \implies (a, b) = 0
    \end{equation*}
    \begin{align*}
        (a, b) \neq (0, 0) \implies \\
        a = b \cdot q_1 + r_1 & \quad 0 < r_1 < |b| \\
        b = r_1 \cdot q_2 + r_2 & \quad 0 < r_2 < r_1 \\
        r_1 = r_2 \cdot q_3 + r_3 & \quad 0 < r_3 < r_2 \\
        \dots \\
        r_{n - 2} = r_{n - 1} \cdot q_n + r_n & \quad 0 < r_n < r_{n - 1} \\
        r_{n - 1} = r_n \cdot q_{n + 1} \\
        (a, b) = r_n
    \end{align*}
\end{proof}
